\section{Children Sum Property}
\label{sec:trees-children-sum}

\problemstatement{The children-sum property holds if, for every node with
at least one child, the node's value equals the sum of its children's
values (a missing child counts as $0$). Given a tree, \emph{modify node
values} (you may only increase them) so the property holds, without
changing the tree's structure.}

\begin{patternbox}
This is not a check but a \textbf{construction}, and the ``you may only
increase values'' rule forces a specific two-pass shape: push required
values \emph{down} on the way in (raise children up to at least the
parent), then fix parents \emph{up} on the way out (set each parent to its
children's actual sum). The down-then-up structure is the key insight.
\end{patternbox}

\subsection*{Understanding the problem carefully}

If we only tried to fix parents to equal their children on the way up, a
parent larger than its children's total would have no legal move (we cannot
decrease). So first, descending, we \emph{raise} the children: give each
child at least the parent's value (bump the larger child, or either, up to
the parent). Now the children's total is $\ge$ the parent, so on the way
back up we can safely set each parent to the exact sum of its (possibly
raised) children.

\begin{ideabox}[Push Down, Then Fix Up]
\textbf{Descent:} at each node, compute the sum of its children's current
values; if that sum is $\ge$ the node, set the node to the sum, else push
the node's value down into (one of) its children. Recurse.
\textbf{Ascent:} after both children are finalised, set the node to the sum
of its children's values. Leaves are left untouched.
\end{ideabox}

\subsection*{Algorithm}

\begin{enumerate}
  \item \textsc{Fix}(\textit{node}):
  \begin{enumerate}
    \item If \textit{node} is null or a leaf, return.
    \item Let \textit{childSum} $=$ (left val if any) $+$ (right val if
          any). If \textit{childSum} $\ge$ \textit{node}.val, set
          \textit{node}.val $=$ \textit{childSum}; else copy
          \textit{node}.val into an existing child.
    \item Recurse into left and right.
    \item Set \textit{node}.val $=$ (left val if any) $+$ (right val if
          any).
  \end{enumerate}
\end{enumerate}

\subsection*{Dry run}

\dryrun{Tree: root $50$; children $7$ and $2$.}

\begin{itemize}
  \item Descent at $50$: childSum $=7+2=9<50$, so push $50$ down: set left
        child to $50$ (now $50,2$).
  \item Recurse into leaves $50$ and $2$: nothing to do.
  \item Ascent at root: set root $=50+2=52$.
  \item Final: root $52$, children $50,2$; $50+2=52$. Property holds.
\end{itemize}

\subsection*{C++ implementation}

\begin{lstlisting}
#include <bits/stdc++.h>
using namespace std;

struct TreeNode {
    int val;
    TreeNode* left;
    TreeNode* right;
    TreeNode(int v) : val(v), left(nullptr), right(nullptr) {}
};

void fix(TreeNode* node) {
    if (node == nullptr) return;
    if (node->left == nullptr && node->right == nullptr) return;

    int childSum = 0;
    if (node->left)  childSum += node->left->val;
    if (node->right) childSum += node->right->val;

    if (childSum >= node->val) {
        node->val = childSum;               // children already suffice
    } else {                                // push node value down
        if (node->left)  node->left->val  = node->val;
        else             node->right->val = node->val;
    }

    fix(node->left);
    fix(node->right);

    int total = 0;                          // fix up from actual children
    if (node->left)  total += node->left->val;
    if (node->right) total += node->right->val;
    node->val = total;
}

int main() {
    TreeNode* root = new TreeNode(50);
    root->left = new TreeNode(7);
    root->right = new TreeNode(2);

    fix(root);
    cout << root->val << " -> "
         << root->left->val << ", " << root->right->val << "\n";
    return 0;
}
\end{lstlisting}

\begin{complexitybox}
\textbf{Time Complexity.} $O(n)$ -- each node is handled once on the way
down and once on the way up. \textbf{Space Complexity.} $O(h)$ for the
recursion stack.
\end{complexitybox}

\begin{takeawaybox}
When a constraint restricts \emph{how} you may change values (here, only
increases), a single pass rarely suffices -- you pre-condition on the way
down so the fix-up on the way back is always legal. Down-then-up is a
recurring shape for tree mutations under monotone constraints.
\end{takeawaybox}
